% DERIVED from formal_layer.json — read-only, do not edit.
\begin{assumptionv}{ass:gaussian}[Gaussian measurement error]
Under \(P\), the ambient latent law, the standardized measurement error \(\varepsilon\) satisfies
\[
\varepsilon\sim N(0,1).
\]
\end{assumptionv}

\begin{assumptionv}{ass:error-independence}[Independent measurement error]
Under \(P\), the measurement error \(\varepsilon\) is independent of the latent running variable \(X\) and the binary potential outcomes \(Y^0,Y^1\):
\[
\varepsilon\perp(X,Y^0,Y^1).
\]
\end{assumptionv}

\begin{definitionv}{synth_2}[Continuous conditional potential-outcome means]
For a latent law \(P\) with score \(X\in[-1,1]\) and binary potential outcomes \(Y^0,Y^1\), let \(\mu_d(P,\cdot)\) be a specified continuous version on \([-1,1]\) of the conditional mean of \(Y^d\) given \(X\), for \(d\in\{0,1\}\). Thus
\[
\mu_d(P,x)=\mathbb E_P[Y^d\mid X=x],
\qquad x\in[-1,1],
\]
with the conditional expectation interpreted through this specified continuous version, including at the endpoints and at the cutoff \(x=0\).
\end{definitionv}

\begin{definitionv}{synth_1}[Latent law and continuous versions]
An ambient latent law \(P\) is a probability law on \(\mathbb R\times\{0,1\}\times\{0,1\}\times\mathbb R\), with coordinates \((X,Y^0,Y^1,\varepsilon)\), together with a density \(f(P,\cdot):\mathbb R\to\mathbb R\) and conditional potential-outcome mean versions \(\mu_d(P,\cdot):\mathbb R\to\mathbb R\), \(d\in\{0,1\}\), satisfying:
\begin{itemize}
\item \textbf{(Latent interval.)} \(P(X\in[-1,1])=1\).
\item \textbf{(Score density.)} The function \(f(P,\cdot)\) is continuous on \([-1,1]\), nonnegative on \(\mathbb R\), and zero outside \([-1,1]\). Moreover,
\[
\int_{-1}^{1} f(P,x)\,dx=1,
\qquad
P(X\in B)=\int_B f(P,x)\,dx
\]
for every measurable set \(B\subseteq\mathbb R\).
\item \textbf{(Continuous mean versions.)} For each \(d\in\{0,1\}\), the function \(\mu_d(P,\cdot)\) is continuous on \([-1,1]\), and for every measurable set \(B\subseteq\mathbb R\),
\[
\int_{\{X\in B\}} Y^d\,dP
=
\int_{B\cap[-1,1]} \mu_d(P,x)f(P,x)\,dx.
\]
\end{itemize}
\end{definitionv}

\begin{assumptionv}{ass:density-bounds}[Latent density bounds]
The latent-score density \(f(P,x)\) satisfies
\[
\frac14\le f(P,x)\le\frac34
\qquad\text{for every }x\in[-1,1].
\]
\end{assumptionv}

\begin{assumptionv}{ass:mean-bounds}[Potential-outcome mean bounds]
The conditional potential-outcome means \(\mu_d(P,x)\) satisfy
\[
\frac14\le\mu_d(P,x)\le\frac34
\qquad\text{for every }d\in\{0,1\}\text{ and }x\in[-1,1].
\]
\end{assumptionv}

\begin{definitionv}{synth_3}[Full-interval Hölder seminorm]
For \(0<\beta\le1\) and a continuous function \(v:[-1,1]\to\mathbb R\), define its full-interval Hölder seminorm by
\[
[v]_\beta
=
\sup_{\substack{x,t\in[-1,1]\\x\ne t}}
\frac{|v(x)-v(t)|}{|x-t|^\beta}.
\]
In particular, for \(d\in\{0,1\}\),
\[
[\mu_d(P,\cdot)]_\beta
=
\sup_{\substack{x,t\in[-1,1]\\x\ne t}}
\frac{|\mu_d(P,x)-\mu_d(P,t)|}{|x-t|^\beta}.
\]
The supremum is taken over the entire latent-score interval and may equal \(+\infty\).
\end{definitionv}

\begin{assumptionv}{ass:mean-holder}[Potential-outcome Hölder bounds]
The full-interval Hölder seminorms \([\mu_d(P,\cdot)]_\beta\) satisfy
\[
[\mu_d(P,\cdot)]_\beta\le2
\qquad\text{for every }d\in\{0,1\}.
\]
\end{assumptionv}

\begin{definitionv}{def:model}[Latent-law class \(\mathcal M_\beta(\sigma)\)]
\(\mathcal M_\beta(\sigma)\) is the class of individual latent laws \(P\) satisfying \cref{obj:ass:gaussian,obj:ass:error-independence,obj:ass:density-bounds,obj:ass:mean-bounds,obj:ass:mean-holder} and the following conditions:
\begin{itemize}
\item \textbf{(Latent support.)} \(P\) is a probability law on
\[
[-1,1]\times\{0,1\}^{2}\times\mathbb R
\]
with coordinates \((X,Y^0,Y^1,\varepsilon)\).
\item \textbf{(Density continuity.)} The latent-score density \(f(P,\cdot)\) is continuous on \([-1,1]\).
\item \textbf{(Density normalization.)}
\[
\int_{-1}^{1}f(P,x)\,dx=1.
\]
\end{itemize}
The Gaussian-contaminated observed score \(W\) is
\[
W=X+\sigma\varepsilon.
\]
For every integer \(n\ge2\), the independent observed sample \(S_n\) has distribution
\[
S_n\sim P_{\mathrm{obs}}^{\otimes n},
\]
where \(P_{\mathrm{obs}}\) denotes the induced observed law.
\end{definitionv}

\begin{definitionv}{def:honest-class}[Honest interval class $\mathcal I_{\beta,n,\sigma}$]
The class of honest connected intervals is
\[
\mathcal I_{\beta,n,\sigma}
=
\left\{
I:(S_n,U)\mapsto[\ell,\upsilon]\subset[-1,1]:
\ell\le\upsilon,\quad
I\text{ is measurable},\quad
\inf_{P\in\mathcal M_\beta(\sigma)}
\mathbb P_{P,n}\{\theta(P)\in I(S_n,U)\}\ge0.9
\right\}.
\]
Here \(S_n\) is the observed sample, \(U\) is the independent procedure seed, \(\mathbb P_{P,n}\) is their joint law under the latent law \(P\), \(\mathcal M_\beta(\sigma)\) is the benchmark law class, and \(\theta(P)\) is the cutoff treatment contrast. The endpoints \(\ell\) and \(\upsilon\) are the outputs of the interval procedure.
\end{definitionv}

\begin{definitionv}{synth_4}[Chebyshev polynomials]
For each nonnegative integer \(L\), let \(T_L\) be the Chebyshev polynomial of the first kind, characterized by
\[
T_L(\cos t)=\cos(Lt),
\qquad t\in\mathbb R.
\]
This normalization gives \(T_0(x)=1\), \(T_1(x)=x\), and \(T_L(1)=1\). Each \(T_L\) is evaluated as a polynomial on all of \(\mathbb R\); the endpoint-kernel construction uses integers \(L\ge1\).
\end{definitionv}

\begin{definitionv}{def:endpoint-kernel}[Endpoint polynomial weight $K_L$]
The endpoint polynomial weight is
\[
K_L(x)
=
\frac{
\left(\dfrac{1-T_L(1-2x)}{2x}\right)^3
}{
\displaystyle\int_0^1
\left(\dfrac{1-T_L(1-2t)}{2t}\right)^3\,dt
}.
\]
Here \(T_L\) is the Chebyshev polynomial. Removable values are filled by continuity, and \(K_L\) is extended polynomially to every real argument.
\end{definitionv}

\begin{definitionv}{def:inverse-heat}[Inverse-heat polynomial $Q_{L,\sigma}$]
The inverse-heat polynomial is
\[
Q_{L,\sigma}(w)
=
\sum_{k=0}^{\lfloor m_L/2\rfloor}
\frac{(-\sigma^2/2)^k}{k!}\,
K_L^{(2k)}(w).
\]
Here \(m_L\) is the degree of the endpoint polynomial \(K_L\), \(\sigma\) is the measurement-error scale, and \(K_L^{(2k)}\) denotes its derivative of order \(2k\).
\end{definitionv}

\begin{definitionv}{synth_5}[Arm-reflection convention]
For the treatment-arm index \(d\in\{0,1\}\), define
\[
s_d=2d-1,
\qquad s_1=1,\qquad s_0=-1.
\]
With \(D=\mathbf1\{X\ge0\}\), the reflected latent score \(s_dX\) lies in \([0,1]\) on the event \(\{D=d\}\). The observable polynomial weight for arm \(d\) is evaluated at the reflected proxy \(s_dW\).
\end{definitionv}

\begin{algorithmv}{def:ratio}[Projected arm estimate $\widehat\mu_{d,L}$]
Inputs are the observed sample \((D_i,Y_i,W_i)_{i=1}^n\), the arm \(d\), the index \(L\), the measurement-error scale \(\sigma\), the inverse-heat polynomial \(Q_{L,\sigma}\), and the arm-reflection sign \(s_d\).
\begin{enumerate}
\item Compute the outcome-marked average:
\[
\widehat A_{d,L}
=
\frac1n\sum_{i=1}^n
\mathbf1\{D_i=d\}Y_iQ_{L,\sigma}(s_dW_i).
\]
\item Compute the unmarked average:
\[
\widehat B_{d,L}
=
\frac1n\sum_{i=1}^n
\mathbf1\{D_i=d\}Q_{L,\sigma}(s_dW_i).
\]
\item Output the projected arm estimate:
\[
\widehat\mu_{d,L}
=
\Pi_{[1/4,3/4]}
\left(
\frac{\widehat A_{d,L}}
{\max\{\widehat B_{d,L},1/8\}}
\right).
\]
Here \(\Pi_{[1/4,3/4]}\) denotes ordinary projection onto \([1/4,3/4]\).
\end{enumerate}
\end{algorithmv}

\begin{definitionv}{synth_6}[Candidate effect estimates, public radii, and intervals]
Fix public parameters \(0<\beta\le1\), an integer \(n\ge2\), and \(0\le\sigma\le1\). For each integer \(L\ge1\), let \(K_L\) be the endpoint kernel, put \(m_L=3(L-1)\), and define the public moment and covariance quantities
\[
M_{L,\beta}=\int_0^1x^\beta K_L(x)\,dx,
\qquad
V_L(\sigma)=\frac34\sum_{j=0}^{m_L}
\frac{\sigma^{2j}}{j!}
\int_0^1[K_L^{(j)}(x)]^2\,dx.
\]
Here \(K_L^{(j)}\) denotes the ordinary \(j\)th derivative, and the \(j=0\) coefficient is one, including at \(\sigma=0\). Using the projected arm estimates, define
\[
\widehat\theta_L=\widehat\mu_{1,L}-\widehat\mu_{0,L},
\qquad
\rho_L=12M_{L,\beta}
+28\sqrt{\frac{40V_L(\sigma)}{n}},
\]
and the connected closed interval
\[
I_L=
[\widehat\theta_L-\rho_L,\widehat\theta_L+\rho_L]
\cap[-1,1].
\]
The index-zero candidate is
\[
\widehat\theta_0=0,\qquad
\rho_0=\frac12,\qquad
I_0=\left[-\frac12,\frac12\right].
\]
Thus each radius is a deterministic function of \((\beta,n,\sigma)\), while each positive-index interval is centered at its observable candidate estimate before intersection with \([-1,1]\).
\end{definitionv}

\begin{algorithmv}{def:public-selection}[Public selection $L_\star$]
Inputs are \((\beta,n,\sigma)\), the finite index cap \(L_{\max}\), and the candidate radii \(\rho_L\), effect estimates \(\widehat\theta_L\), and intervals \(I_L\); index zero uses the specified constant estimate and the full target-range interval.
\begin{enumerate}
\item Select the smallest index minimizing the radius:
\[
L_\star
=
\min\operatorname*{arg\,min}_{L\in\{0,\ldots,L_{\max}\}}
\rho_L.
\]
All comparisons are deterministic functions of \((\beta,n,\sigma)\).
\item Define the selected radius:
\[
E_\beta(n,\sigma)=\rho_{L_\star}.
\]
\item Output the selected estimate and interval:
\[
\bigl(\widehat\theta_{L_\star},I_{L_\star}\bigr).
\]
\end{enumerate}
\end{algorithmv}

\begin{theoremv}{thm:finite-certificate}[Finite sample certificate]
Suppose that:
\begin{itemize}
\item \textbf{(Smoothness.)} The public smoothness parameter satisfies \(0<\beta\le1\).
\item \textbf{(Noise.)} The noise scale satisfies \(0\le\sigma\le1\).
\item \textbf{(Sample and degree.)} The sample size \(n\) and polynomial index \(L\) are integers satisfying \(n\ge2\) and \(L\ge1\).
\item \textbf{(Latent law.)} \(P\in\mathcal M_\beta(\sigma)\).
\end{itemize}
Write
\[
\theta(P)=\mu_1(P,0)-\mu_0(P,0),
\qquad
B_{d,L}(P)=\mathbb E_{P,n}\widehat B_{d,L},
\]
where \(\mathbb P_{P,n}\) is the joint law of the independent observed sample and independent uniform procedure seed, and \(\mathbb E_{P,n}\) denotes expectation under this law. Define the endpoint bias moment and derivative variance cost by
\[
M_{L,\beta}
=\int_0^1 x^\beta K_L(x)\,dx,
\qquad
V_L(\sigma)
=\frac34\sum_{j=0}^{m_L}
\frac{\sigma^{2j}}{j!}
\int_0^1\bigl(K_L^{(j)}(x)\bigr)^2\,dx,
\]
where \(m_L\) is the degree of the endpoint polynomial \(K_L\).
For the effect estimate
\[
\widehat\theta_L=\widehat\mu_{1,L}-\widehat\mu_{0,L}
\]
and its connected closed interval \(I_L\), the following guarantees hold:
\[
B_{d,L}(P)\ge\frac14
\quad\text{for every }d\in\{0,1\},
\]
\[
\mathbb E_{P,n}\bigl|\widehat\theta_L-\theta(P)\bigr|
\le
12M_{L,\beta}
+28\sqrt{\frac{V_L(\sigma)}{n}},
\qquad
\mathbb P_{P,n}\{\theta(P)\in I_L\}\ge\frac9{10}.
\]
Moreover, the publicly selected interval \(I_{L_\star}\) belongs to the uniformly honest class \(\mathcal I_{\beta,n,\sigma}\) of \cref{obj:def:honest-class}, and the selected estimator and interval satisfy
\[
\sup_{P\in\mathcal M_\beta(\sigma)}
\mathbb E_{P,n}
\bigl|\widehat\theta_{L_\star}-\theta(P)\bigr|
\le E_\beta(n,\sigma),
\qquad
\sup_{P\in\mathcal M_\beta(\sigma)}
\mathbb E_{P,n}\operatorname{length}(I_{L_\star})
\le2E_\beta(n,\sigma).
\]
\end{theoremv}

\begin{definitionv}{def:uniform-rate}[Uniform rate \(r_\beta(n,\sigma)\)]
The uniform resolution rate and selected observable procedures are
\[
r_\beta(n,\sigma)=h_\star(n,\sigma)^\beta,
\qquad
\widetilde\theta_{\beta,n,\sigma}=\widehat\theta_{L_\star},
\qquad
\widetilde I_{\beta,n,\sigma}=I_{L_\star},
\]
where \(L_\star\) is selected by the finite public minimization. With the local abbreviations
\[
\nu=2\beta+1,\qquad h_0=n^{-1/\nu},
\]
define the noise cost, for \(0<h\le1\), by
\[
\Psi(h,\sigma)=
\begin{cases}
0,&\sigma\le h,\\
(\sigma/h)^{2/3}-1,&\sigma^4\le h<\sigma,\\
h^{-1/2}\{1+\log(\sigma^2/\sqrt h)\}-1,&h<\sigma^4.
\end{cases}
\]
At \(\sigma=0\), the first branch applies throughout \(0<h\le1\). The resolution \(h_\star(n,\sigma)\) is specified as the root in \([h_0,1]\) of
\[
\log n+\nu\log h_\star(n,\sigma)
=\Psi(h_\star(n,\sigma),\sigma).
\]
The left side is continuous and increasing in the resolution, and the right side is continuous and nonincreasing. Root existence and uniqueness, risk bounds, and interval honesty are subsequent assertions. All tuning parameters belong to the procedures; the law class retains its defining restrictions.
\end{definitionv}

\begin{theoremv}{thm:uniform-frontier}[Uniform estimation and interval frontier]
Fix \(\beta\in(0,1]\). There exist constants \(c,C>0\), depending only on \(\beta\), such that the following conclusions hold simultaneously for every \(n\ge2\) and every \(\sigma\in[0,1]\).

Write \(\nu=2\beta+1\) and \(h_0=n^{-1/\nu}\). The equation defining the resolution \(h_\star(n,\sigma)\) has a unique root. With \(r_\beta(n,\sigma)=h_\star(n,\sigma)^\beta\), the selected observable estimator \(\widetilde\theta_{\beta,n,\sigma}\), the selected observable interval \(\widetilde I_{\beta,n,\sigma}\), and their public radius \(E_\beta(n,\sigma)\) satisfy
\[
c r_\beta(n,\sigma)
\le R_\beta(n,\sigma)
\le \sup_{P\in\mathcal M_\beta(\sigma)}
\mathbb E_{P,n}
\left|\widetilde\theta_{\beta,n,\sigma}-\theta(P)\right|
\le E_\beta(n,\sigma)
\le C r_\beta(n,\sigma)
\]
and
\[
c r_\beta(n,\sigma)
\le \mathcal L_\beta(n,\sigma)
\le \sup_{P\in\mathcal M_\beta(\sigma)}
\mathbb E_{P,n}
\operatorname{length}(\widetilde I_{\beta,n,\sigma})
\le 2E_\beta(n,\sigma)
\le C r_\beta(n,\sigma),
\qquad
\widetilde I_{\beta,n,\sigma}\in\mathcal I_{\beta,n,\sigma}.
\]
The minimax criteria range over all procedures in their stated decision classes, including procedures with arbitrary independent randomization and arbitrary degree.

The following branches exhaust every sample size. Throughout, \(\asymp\) denotes comparison by positive constants depending only on \(\beta\), except where dependence on an additional fixed exponent is stated.
\begin{itemize}
\item \textbf{(Direct branch.)}
If \(\sigma\le h_0\), then \(h_\star=h_0\), including at \(\sigma=0\).

\item \textbf{(Intermediate branch.)}
If \(\sigma>h_0\), then \(h_0<h_\star<\sigma\). Within this noisy region, the intermediate branch occurs exactly when
\[
\log(n\sigma^{4\nu})\le \sigma^{-2}-1.
\]
There is a unique \(z\in(1,\sigma^{-2}]\) satisfying
\[
z+\frac{3\nu}{2}\log z=1+\log(n\sigma^\nu),
\]
and
\[
h_\star=\sigma z^{-3/2}
\asymp
\frac{\sigma}{[1+\log(n\sigma^\nu)]^{3/2}}.
\]

\item \textbf{(Compact branch.)}
If \(\sigma>h_0\) and
\[
\log(n\sigma^{4\nu})>\sigma^{-2}-1,
\]
there is a unique \(\tau>\sigma^{-2}\) satisfying
\[
2\nu\log\tau+
\tau\{1+\log(\sigma^2\tau)\}
=\log n+1.
\]
On this branch,
\[
h_\star=\tau^{-2},
\qquad
\tau\asymp
\frac{\log(en)}{\log(e+\sigma^2\log(en))},
\qquad
r_\beta(n,\sigma)\asymp
\left\{
\frac{\log(e+\sigma^2\log(en))}{\log(en)}
\right\}^{2\beta},
\]
uniformly over the branch.
\end{itemize}
At the intermediate--compact equality, \(h_\star=\sigma^4\), and both scalar formulas give this value. At the direct--noisy equality, \(h_\star=\sigma=h_0\), and the adjacent formulas agree.

For every fixed \(a\) with \(0<a<1/(2\beta+1)\),
\[
r_\beta(n,n^{-a})
\asymp n^{-a\beta}(\log n)^{-3\beta/2}
\]
as \(n\) tends to infinity, with comparison constants that may additionally depend on \(a\). For every fixed \(a\ge1/(2\beta+1)\), the sequence \(\sigma=n^{-a}\) lies in the exact direct branch for every \(n\ge2\). For every fixed \(\sigma\in(0,1]\),
\[
r_\beta(n,\sigma)
\asymp
\left\{\frac{\log\log(en)}{\log(en)}\right\}^{2\beta}
\]
as \(n\) tends to infinity, with comparison constants allowed to depend on the fixed noise level. The root and finite-sample branch formulas characterize every shrinking-noise sequence.

In particular, for \(\beta=1\) and \(\sigma=n^{-1/6}\), the intermediate branch holds for every \(n\ge2\), and
\[
r_1(n,n^{-1/6})
=h_\star(n,n^{-1/6})
\asymp
\frac{n^{-1/6}}{[1+\tfrac12\log n]^{3/2}}.
\]
\end{theoremv}

\begin{propositionv}{prop:published-class-converse-transfer}[Gaussian class converse transfer]
For \(\sigma\in[0,1]\), let \(\mathcal G(\sigma)\) consist of probability laws \(Q\) on \(\mathbb R\times\{0,1\}^2\times\mathbb R\), with coordinates \((Z,V^0,V^1,E)\), a score density \(a\), and conditional-mean versions
\[
v_d(z)=\mathbb E_Q[V^d\mid Z=z],\qquad d\in\{0,1\}.
\]
Define the assignment, observed outcome, and observed score by
\[
A=\mathbf 1\{Z\ge0\},\qquad
B=AV^1+(1-A)V^0,\qquad
R=Z+\sigma E.
\]
Membership requires the following conditions:
\begin{itemize}
\item \textbf{(Local density.)} There exists \(r>0\) such that \(a\) is continuous on \((-r,r)\) and \(a(z)>0\) for every \(z\in(-r,r)\).
\item \textbf{(Conditional means.)} Each \(v_d\) is a conditional-mean version continuous on an open neighborhood of zero, with \(0\le v_d(0)\le1\).
\item \textbf{(Gaussian error.)} \(E\sim N(0,1)\).
\item \textbf{(Joint independence.)} \(E\perp(Z,B)\).
\end{itemize}
The target and randomized experiment are
\[
\vartheta(Q)=v_1(0)-v_0(0),\qquad
\bigl(Q\circ(R,A,B)^{-1}\bigr)^{\otimes n}
 \otimes\operatorname{Unif}[0,1].
\]
Let \(\mathcal R_{\mathcal G}(n,\sigma)\) denote the infimum, over measurable estimates based on this experiment, of
\[
\sup_{Q\in\mathcal G(\sigma)}
\mathbb E_Q\!\left[\,|T-\vartheta(Q)|\,\right].
\]
Let \(\mathcal L_{\mathcal G}(n,\sigma)\) denote the infimum of
\[
\sup_{Q\in\mathcal G(\sigma)}
\mathbb E_Q[\upsilon-\ell]
\]
over measurable connected closed intervals \([\ell,\upsilon]\subseteq[-1,1]\) based on the same experiment and satisfying
\[
\inf_{Q\in\mathcal G(\sigma)}
Q\{\vartheta(Q)\in[\ell,\upsilon]\}\ge0.9.
\]
These expectations and coverage probabilities include the independent uniform procedure seed.

Fix \(\beta\in(0,1]\). For every \(\sigma\in[0,1]\) and every \(P\in\mathcal M_\beta(\sigma)\), the law \(P\), equipped with the density \(a=f(P,\cdot)\) and versions \(v_d=\mu_d(P,\cdot)\), belongs to \(\mathcal G(\sigma)\). This embedding preserves the target,
\[
\vartheta(Q)=\theta(P)
=\mu_1(P,0)-\mu_0(P,0),
\]
the observed law,
\[
Q\circ(R,A,B)^{-1}=P_{\mathrm{obs}},
\]
and, for every integer \(n\ge0\), the entire randomized experiment,
\[
\bigl(Q\circ(R,A,B)^{-1}\bigr)^{\otimes n}
 \otimes\operatorname{Unif}[0,1]
=\mathbb P_{P,n}.
\]

There exists \(c>0\), depending on \(\beta\), such that for every integer \(n\ge2\) and every \(\sigma\in[0,1]\),
\[
c\,r_\beta(n,\sigma)
\le R_\beta(n,\sigma)
\le \mathcal R_{\mathcal G}(n,\sigma),
\qquad
c\,r_\beta(n,\sigma)
\le \mathcal L_\beta(n,\sigma)
\le \mathcal L_{\mathcal G}(n,\sigma),
\]
where \(r_\beta(n,\sigma)\) is the benchmark frontier rate, and \(R_\beta(n,\sigma)\) and \(\mathcal L_\beta(n,\sigma)\) are the corresponding minimax absolute risk and minimax expected length under uniform \(0.9\) coverage on \(\mathcal M_\beta(\sigma)\).

For every \(\sigma\in(0,1]\) and every \(Q\in\mathcal G(\sigma)\), the locally defined binary-response Pei--Shen conditions hold: the standardized error is independent jointly of the latent score and observed outcome (Assumption 1Y), the score density is strictly positive on an open neighborhood of zero (Assumption 2C), and the scaled error satisfies
\[
\sigma E\sim N(0,\sigma^2)
\]
(Assumption 7).
\end{propositionv}

\begin{theoremv}{lem:classical-hermite-facts}[Gaussian Hermite orthogonality]
Let \(H_j\) be the probabilists' Hermite polynomials, defined by
\[
\exp(zt-t^2/2)
=\sum_{j=0}^{\infty}H_j(z)\frac{t^j}{j!}.
\]
For a standard Gaussian variable \(Z\) and nonnegative integers \(j,k\),
\[
\mathbb E[H_j(Z)H_k(Z)]
=j!\mathbf1\{j=k\},
\qquad H_0=1.
\]
\end{theoremv}

\begin{lemmav}{lem:positive-endpoint}[Positive endpoint kernel bounds]
Let \(0<\beta\le1\), and let \(K_L\) be the endpoint kernel defined in \cref{obj:def:endpoint-kernel}. Define its bias moment by
\[
M_{L,\beta}=\int_0^1 x^\beta K_L(x)\,dx.
\]
There exists a constant \(C>0\), depending on \(\beta\), such that for every integer \(L\ge1\),
\[
K_L(x)\ge0\quad\text{for all }x\in[0,1],
\qquad
\int_0^1 K_L(x)\,dx=1,
\]
and
\[
M_{L,\beta}\le C L^{-2\beta},
\qquad
\int_0^1 K_L(x)^2\,dx\le C L^2.
\]
\end{lemmav}

\begin{lemmav}{lem:factorial-derivatives}[Factorial polynomial derivative bound]
Let \(m,j\) be nonnegative integers with \(1\le j\le m\), and let \(p\) be a real polynomial of degree at most \(m\). Writing \(p^{(j)}\) for its \(j\)-th ordinary derivative, we have
\[
\left(\int_0^1 \bigl(p^{(j)}(x)\bigr)^2\,dx\right)^{1/2}
\le
\frac{4^j(m+1)^{2j}}{j!}
\left(\int_0^1 p(x)^2\,dx\right)^{1/2}.
\]
\end{lemmav}

\begin{lemmav}{lem:exact-covariance}[Exact Gaussian moment identities]
There exists an absolute constant \(C>0\) such that, for every integer \(L\ge1\) and every \(\sigma\in[0,1]\), the following holds. Let \(K_L\) and \(Q_{L,\sigma}\) be defined by \cref{obj:def:endpoint-kernel,obj:def:inverse-heat}, and set
\[
m_L=3(L-1),\qquad
V_L(\sigma)=\frac34\sum_{j=0}^{m_L}\frac{\sigma^{2j}}{j!}
\int_0^1 [K_L^{(j)}(x)]^2\,dx.
\]
For a standard Gaussian random variable \(\varepsilon\), at every \(x\in\mathbb R\),
\[
\mathbb E\!\left[Q_{L,\sigma}(x+\sigma\varepsilon)\right]=K_L(x),
\qquad
\mathbb E\!\left[Q_{L,\sigma}(x+\sigma\varepsilon)^2\right]
=\sum_{j=0}^{m_L}\frac{\sigma^{2j}[K_L^{(j)}(x)]^2}{j!}.
\]
Moreover,
\[
V_L(\sigma)\le C L^2
\sum_{j=0}^{m_L}\frac{(C\sigma^2L^4)^j}{(j!)^3}.
\]
\end{lemmav}

\begin{theoremv}{lem:classical-legendre-facts}[Classical Legendre polynomial identities]
For the normalized Legendre polynomials \(P_j^{\mathrm{Leg}}\), every pair of nonnegative integers \(j,k\) satisfies
\[
\int_{-1}^{1} P_j^{\mathrm{Leg}}(t)P_k^{\mathrm{Leg}}(t)\,dt
=
\begin{cases}
\dfrac{2}{2j+1}, & j=k,\\
0, & j\ne k.
\end{cases}
\]
For every nonnegative integer \(j\), reflection symmetry holds for every \(t\in\mathbb R\):
\[
P_j^{\mathrm{Leg}}(-t)=(-1)^jP_j^{\mathrm{Leg}}(t).
\]
The same polynomial satisfies
\[
|P_j^{\mathrm{Leg}}(t)|\le 1
\quad\text{for every }t\in[-1,1],
\qquad
P_j^{\mathrm{Leg}}(-1)=(-1)^j.
\]
\end{theoremv}

\begin{lemmav}{lem:shifted-legendre-derivative}[Shifted Legendre derivative bounds]
Let \(P_k^{\mathrm{Leg}}\) denote the degree-\(k\) Legendre polynomial normalized by \(P_k^{\mathrm{Leg}}(1)=1\), and define
\[
R_k(t)=P_k^{\mathrm{Leg}}(2t-1),\qquad k\in\mathbb N,\quad t\in\mathbb R.
\]
For every integer \(j\ge1\),
\[
R_j'(t)=2\sum_{\substack{0\le k<j\\j-k\ \mathrm{odd}}}(2k+1)R_k(t)
\qquad\text{for all }t\in\mathbb R,
\]
and
\[
|R_j'(t)|\le j(j+1)
\qquad\text{for all }t\in[0,1].
\]
\end{lemmav}

\begin{definitionv}{def:legal-extension}[Cancellation block and full-interval extension $v_{b,m}$]
For an integer \(m\ge2\), define the block normalizer, cancellation polynomial, and fixed amplitude by
\[
N_m=(2m+1)^2-m^2,
\qquad
g_m(t)=\frac{1-t}{N_m}
\sum_{j=m}^{2m}(2j+1)(-1)^jP_j^{\mathrm{Leg}}(2t-1),
\qquad
\kappa=\frac1{100}.
\]
For \(0<b\le1\), the full-interval extension is
\[
v_{b,m}(x)
=
\begin{cases}
1, & -1\le x<0,\\
g_m(x/b), & 0\le x\le b,\\
0, & b<x\le1.
\end{cases}
\]
\end{definitionv}

\begin{definitionv}{def:alternatives}[Alternative laws \(P^\pm_{b,m}\)]
The latent laws \(P^\pm_{b,m}\) are specified by:
\begin{itemize}
\item \textbf{(Latent score.)}
\[
X\sim\operatorname{Uniform}[-1,1].
\]
\item \textbf{(Potential outcomes.)} Conditional on \(X=x\), the binary potential outcomes \(Y^0\) and \(Y^1\) are independent, with
\[
Y^0\mid X=x\sim\operatorname{Bernoulli}(1/2),
\qquad
Y^1\mid X=x\sim
\operatorname{Bernoulli}\!\left(
\frac12\pm\kappa(b/m^2)^\beta v_{b,m}(x)
\right).
\]
\item \textbf{(Measurement error.)} The standardized measurement error \(\varepsilon\) is independent of \((X,Y^0,Y^1)\), with
\[
\varepsilon\sim N(0,1).
\]
\end{itemize}
The induced observed law uses the benchmark observation map
\[
W=X+\sigma\varepsilon,\qquad
D=\mathbf1\{X\ge0\},\qquad
Y=(1-D)Y^0+DY^1.
\]
\end{definitionv}

\begin{lemmav}{lem:legal-cancellation}[Legal cancellation and separation]
Let \(\sigma\in\mathbb R\), and suppose that:
\begin{itemize}
\item \textbf{(Smoothness.)} \(0<\beta\le1\).
\item \textbf{(Support width.)} \(0<b\le1\).
\item \textbf{(Polynomial order.)} \(m\in\mathbb N\) and \(m\ge2\).
\end{itemize}
For the normalizer \(N_m\), cancellation polynomial \(g_m\), amplitude \(\kappa\), and extension \(v_{b,m}\) defined in \cref{obj:def:legal-extension},
\[
g_m(0)=1,\qquad g_m(1)=0,\qquad
\sup_{t\in[0,1]}|g_m(t)|\le1,\qquad
\int_0^1g_m(t)^2\,dt\le\frac1{N_m},
\]
and, for every integer \(k\) with \(0\le k\le m-2\),
\[
\int_0^1t^kg_m(t)\,dt=0.
\]
Moreover,
\[
\sup_{t\in[0,1]}|g_m^{(1)}(t)|\le7m^2.
\]
For the extension \(v_{b,m}\) in \cref{obj:def:legal-extension}, its full-interval Hölder seminorm satisfies
\[
[v_{b,m}]_\beta
=
\sup_{\substack{x,t\in[-1,1]\\x\ne t}}
\frac{|v_{b,m}(x)-v_{b,m}(t)|}{|x-t|^\beta}
\le7\left(\frac{b}{m^2}\right)^{-\beta}.
\]
The two latent alternatives \(P^\pm_{b,m}\) of \cref{obj:def:alternatives} both belong to \(\mathcal M_\beta(\sigma)\). Their cutoff causal contrasts, where \(\theta(P)\) is the treated conditional potential-outcome mean at zero minus the untreated conditional potential-outcome mean at zero, satisfy
\[
\left|\theta(P^+_{b,m})-\theta(P^-_{b,m})\right|
=
2\kappa\left(\frac{b}{m^2}\right)^\beta.
\]
\end{lemmav}

\begin{lemmav}{lem:full-marked-likelihood}[Full marked likelihood bound]
Consider the alternative laws \(P^\pm_{b,m}\) of \cref{obj:def:alternatives}, with the cancellation polynomial \(g_m\) of \cref{obj:lem:legal-cancellation}, under the following parameter conditions:
\begin{itemize}
\item \textbf{(Smoothness.)} \(0<\beta\le1\).
\item \textbf{(Support.)} \(0<b\le1\).
\item \textbf{(Order.)} \(m\) is an integer with \(m\ge2\).
\item \textbf{(Noise.)} \(0<\sigma\le1\).
\end{itemize}
Set
\[
\kappa=\frac1{100},
\qquad
N_m=(2m+1)^2-m^2,
\qquad
\lambda_{b,\sigma}=\frac{b^2}{4\sigma^2},
\]
and define the likelihood-comparison quantities by
\[
q_\sigma(w)
=\int_0^1
\frac{\exp\!\bigl(-(w-x)^2/(2\sigma^2)\bigr)}
{\sqrt{2\pi}\,\sigma}\,dx,
\qquad
u_{b,m,\sigma}(w)
=\int_0^b g_m(x/b)
\frac{\exp\!\bigl(-(w-x)^2/(2\sigma^2)\bigr)}
{\sqrt{2\pi}\,\sigma}\,dx.
\]
Here \(\chi^2\) denotes the extended chi-squared divergence: for probability laws, it is the integral of the squared Radon--Nikodym derivative minus one when the first law is absolutely continuous with respect to the second, and \(+\infty\) otherwise.

The full observed laws satisfy
\[
\begin{aligned}
&\chi^2\!\left(
(P^+_{b,m})_{\mathrm{obs}},
(P^-_{b,m})_{\mathrm{obs}}
\right)\\
&\quad\le
16\kappa^2\left(\frac{b}{m^2}\right)^{2\beta}
\int_{\mathbb R}
\frac{u_{b,m,\sigma}(w)^2}{q_\sigma(w)}\,dw\\
&\quad\le
16\kappa^2\left(\frac{b}{m^2}\right)^{2\beta}
\frac{b}{N_m}
\exp\!\left(\frac{b^2}{24\sigma^2}\right)
\sum_{j=m-1}^{\infty}\frac{\lambda_{b,\sigma}^{\,j}}{j!}.
\end{aligned}
\]
Moreover,
\[
(P^+_{b,m})_{\mathrm{obs}}
\ll
(P^-_{b,m})_{\mathrm{obs}}.
\]
The squared likelihood-ratio deviation
\[
\left(
\frac{d(P^+_{b,m})_{\mathrm{obs}}}
{d(P^-_{b,m})_{\mathrm{obs}}}-1
\right)^2
\]
is integrable under \((P^-_{b,m})_{\mathrm{obs}}\), the function
\(w\mapsto u_{b,m,\sigma}(w)^2/q_\sigma(w)\) is Lebesgue integrable, and the displayed factorial series converges. The divergence compares the full observed triples, including both binary outcome cells and all real proxy values.
\end{lemmav}

\begin{lemmav}{lem:two-point-transfer}[Two-point minimax risk transfer]
Let \(P^+_{b,m}\) and \(P^-_{b,m}\) be the alternative laws from \cref{obj:def:alternatives}, constructed using \cref{obj:lem:classical-legendre-facts}. Suppose:
\begin{itemize}
\item \textbf{(Parameter ranges.)} \(0<\beta\le1\), \(0<b\le1\), and \(0\le\sigma\le1\).
\item \textbf{(Sample size and degree.)} \(n,m\in\mathbb N\) with \(n\ge2\) and \(m\ge2\).
\item \textbf{(Sample-law proximity.)} The independent observed-sample laws satisfy
\[
\operatorname{TV}\!\left(
(P^+_{b,m})_{\mathrm{obs}}^{\otimes n},
(P^-_{b,m})_{\mathrm{obs}}^{\otimes n}
\right)\le\frac15,
\]
where total variation is the supremum, over measurable events, of the absolute difference between their probabilities.
\end{itemize}
Set \(\kappa=1/100\). Write \(R_\beta(n,\sigma)\) for the benchmark minimax expected absolute error in the cutoff treatment contrast, and \(\mathcal L_\beta(n,\sigma)\) for the benchmark minimax worst expected length of measurable connected closed intervals with uniform \(0.9\) coverage. Both criteria minimize extended nonnegative losses before converting the resulting minimax values to real numbers; their decision classes allow independent procedure randomization. Then
\[
R_\beta(n,\sigma)\ge
\frac45\,\kappa\left(\frac{b}{m^2}\right)^\beta
\qquad\text{and}\qquad
\mathcal L_\beta(n,\sigma)\ge
\frac65\,\kappa\left(\frac{b}{m^2}\right)^\beta.
\]
\end{lemmav}

\begin{definitionv}{synth_7}[Direct-experiment witness profile]
For \(0<h\le1\), define the direct-experiment witness profile on the full latent-score interval by
\[
v_h^{\mathrm{dir}}(x)=
\begin{cases}
1,&-1\le x<0,\\
\max\{1-x/h,0\},&0\le x\le1.
\end{cases}
\]
This profile is continuous on \([-1,1]\), takes values in \([0,1]\), satisfies \(v_h^{\mathrm{dir}}(0)=1\), and vanishes for \(h\le x\le1\). Its full-interval regularity satisfies
\[
|v_h^{\mathrm{dir}}(x)-v_h^{\mathrm{dir}}(t)|
\le\min\{1,|x-t|/h\},
\qquad x,t\in[-1,1],
\]
and hence, for every \(0<\beta\le1\),
\[
[v_h^{\mathrm{dir}}]_\beta\le h^{-\beta}.
\]
The direct latent witnesses use this profile through the conditional means
\[
\mu_0(P_h^{\mathrm{dir},\pm},x)=\frac12,
\qquad
\mu_1(P_h^{\mathrm{dir},\pm},x)
=\frac12\pm\kappa h^\beta v_h^{\mathrm{dir}}(x),
\qquad \kappa=\frac1{100}.
\]
\end{definitionv}

\begin{definitionv}{def:direct-alternatives}[Direct alternative laws \(P_h^{\mathrm{dir},\pm}\)]
The latent laws \(P_h^{\mathrm{dir},\pm}\) are specified by:
\begin{itemize}
\item \textbf{(Latent score.)}
\[
X\sim\operatorname{Uniform}[-1,1].
\]
\item \textbf{(Potential outcomes.)} Conditional on \(X=x\), the binary potential outcomes are independent Bernoulli variables with conditional means
\[
\mu_0(P_h^{\mathrm{dir},\pm},x)=\frac12,
\qquad
\mu_1(P_h^{\mathrm{dir},\pm},x)
=\frac12\pm\kappa h^\beta v_h^{\mathrm{dir}}(x),
\]
where \(v_h^{\mathrm{dir}}\) is the direct-experiment profile.
\item \textbf{(Measurement error.)} The standardized measurement error \(\varepsilon\) is independent of the latent score and potential outcomes, with
\[
\varepsilon\sim N(0,1).
\]
\end{itemize}
\end{definitionv}

\begin{definitionv}{def:noise-support}[Support and resolution \(b_m(\sigma),h_m(\sigma)\)]
The noise-dependent support scale \(b_m(\sigma)\) and endpoint resolution \(h_m(\sigma)\) are
\[
b_m(\sigma)=\min\{1,\sigma\sqrt m/8\},
\qquad
h_m(\sigma)=\frac{b_m(\sigma)}{m^2},
\qquad \sigma>0.
\]
The corresponding alternative pair is
\[
P^\pm_{b_m(\sigma),m}.
\]
\end{definitionv}

\begin{lemmav}{lem:noise-cost-scaling}[Noise cost scaling bounds]
For \(h\in(0,1]\) and \(\sigma\in[0,1]\), define the noise cost
\[
\Psi(h,\sigma)=
\begin{cases}
0, & \sigma\le h,\\
(\sigma/h)^{2/3}-1, & h<\sigma\ \text{and}\ \sigma^4\le h,\\
h^{-1/2}\bigl[1+\log(\sigma^2/\sqrt h)\bigr]-1,
& h<\sigma\ \text{and}\ h<\sigma^4.
\end{cases}
\]
For every \(\sigma\in[0,1]\), the function \(h\mapsto\Psi(h,\sigma)\) is continuous, nonnegative, and nonincreasing on \((0,1]\).

For all real \(h,t,\sigma\) satisfying \(0<h\le t\le\sigma\le1\),
\[
(t/h)^{1/2}
\le
\frac{1+\Psi(h,\sigma)}{1+\Psi(t,\sigma)}
\le t/h.
\]
For every \(h\in(0,1]\), \(\sigma\in[0,1]\), and real \(a\ge1\) satisfying \(ah\le1\),
\[
1+\Psi(ah,\sigma)
\le
\max\bigl\{1,\ a^{-1/2}[1+\Psi(h,\sigma)]\bigr\}.
\]
\end{lemmav}

\begin{lemmav}{lem:variance-cost-envelope}[Uniform variance cost envelope]
Let \(m_L\) be the degree of the endpoint polynomial \(K_L\), and define its integrated derivative cost by
\[
V_L(\sigma)
=\frac34\sum_{j=0}^{m_L}\frac{\sigma^{2j}}{j!}
\int_0^1\bigl(K_L^{(j)}(x)\bigr)^2\,dx.
\]
For \(h>0\), define the noise cost by
\[
\Psi(h,\sigma)=
\begin{cases}
0, & \sigma\le h,\\
(\sigma/h)^{2/3}-1, & \sigma>h\ \text{and}\ \sigma^4\le h,\\
h^{-1/2}\bigl[1+\log(\sigma^2/\sqrt h)\bigr]-1,
& \sigma>h\ \text{and}\ \sigma^4>h.
\end{cases}
\]
There exists an absolute constant \(C>0\) such that, for every integer \(L\ge1\) and every \(\sigma\in[0,1]\),
\[
V_L(\sigma)\le C L^2
\exp\!\left\{C\bigl[1+\Psi(L^{-2},\sigma)\bigr]\right\}.
\]
\end{lemmav}

\begin{lemmav}{lem:cancellation-information-envelope}[Cancellation information envelope]
Let the parameters satisfy:
\begin{itemize}
\item \textbf{(Smoothness.)} \(0<\beta\le1\).
\item \textbf{(Noise scale.)} \(0<\sigma\le1\).
\item \textbf{(Polynomial degree.)} \(m\in\mathbb N\) and \(m\ge2\).
\end{itemize}
Write \(b=b_m(\sigma)\) and \(h=h_m(\sigma)\) as in
\cref{obj:def:noise-support}, and let \(P^+_{b,m}\) and \(P^-_{b,m}\)
be the alternative laws of \cref{obj:def:alternatives}, constructed using
\cref{obj:lem:classical-legendre-facts}. Set the perturbation constant
\(\kappa=1/100\), and define the noise cost by
\[
\Psi(h,\sigma)=
\begin{cases}
0, & \sigma\le h,\\
(\sigma/h)^{2/3}-1, & \sigma>h\ \text{and}\ \sigma^4\le h,\\
h^{-1/2}\bigl(1+\log(\sigma^2/\sqrt h)\bigr)-1,
& \sigma>h\ \text{and}\ \sigma^4>h.
\end{cases}
\]
For probability laws \(\mu,\nu\), use the extended chi-squared divergence
\[
\chi^2(\mu,\nu)=
\begin{cases}
\displaystyle\int\left(\frac{d\mu}{d\nu}-1\right)^2\,d\nu,
& \mu\ll\nu,\\
+\infty, & \mu\not\ll\nu.
\end{cases}
\]
Then the observed laws induced by the two alternatives at noise scale
\(\sigma\) satisfy
\[
\chi^2\bigl((P^+_{b,m})_{\mathrm{obs}},
           (P^-_{b,m})_{\mathrm{obs}}\bigr)
\le
\frac{32}{3}\kappa^2 h^{2\beta+1}
\exp\!\left(-\frac{1+\Psi(h,\sigma)}{32}\right).
\]
\end{lemmav}

\begin{theoremv}{thm:direct-reduction}[Direct-experiment rates and attainment]
Fix \(0<\beta\le1\), where \(\beta\) is the Hölder exponent. Write \(R_\beta(n,\sigma)\) and \(\mathcal L_\beta(n,\sigma)\) for the real conversions of the extended minimax absolute risk and extended minimax honest expected length, respectively. Their extended objectives are
\[
\inf_{\widehat\theta}\sup_{P\in\mathcal M_\beta(\sigma)}
\mathbb E_{P,n}\!\left[|\widehat\theta-\theta(P)|\right],
\qquad
\inf_{I\in\mathcal I_{\beta,n,\sigma}}
\sup_{P\in\mathcal M_\beta(\sigma)}
\mathbb E_{P,n}\!\left[\operatorname{length}(I)\right].
\]
Here the first infimum ranges over admissible estimators, \(\mathcal M_\beta(\sigma)\) is the benchmark model class, \(\theta(P)\) is the cutoff treatment contrast, and \(\mathbb E_{P,n}\) denotes expectation under the joint law of the observed sample and independent procedure seed. The honest interval class \(\mathcal I_{\beta,n,\sigma}\) is defined in \cref{obj:def:honest-class}. Both minimizations are performed in the extended nonnegative reals before conversion to real numbers.

There exist constants \(c,C>0\), depending only on \(\beta\), such that, for every integer sample size \(n\ge2\) and every noise scale \(\sigma\in[0,1]\),
\[
R_\beta(n,\sigma)\ge c\,n^{-\beta/(2\beta+1)},
\qquad
\mathcal L_\beta(n,\sigma)\ge c\,n^{-\beta/(2\beta+1)}.
\]
At \(\sigma=0\),
\[
R_\beta(n,0)\le C\,n^{-\beta/(2\beta+1)},
\qquad
\mathcal L_\beta(n,0)\le C\,n^{-\beta/(2\beta+1)}.
\]
The selected estimator \(\widehat\theta_{L_\star}\) and selected interval \(I_{L_\star}\), evaluated at \((\beta,n,0)\), satisfy the extended worst-case bounds
\[
\sup_{P\in\mathcal M_\beta(0)}
\mathbb E_{P,n}\!\left[
|\widehat\theta_{L_\star}-\theta(P)|
\right]
\le C\,n^{-\beta/(2\beta+1)},
\qquad
\sup_{P\in\mathcal M_\beta(0)}
\mathbb E_{P,n}\!\left[
\operatorname{length}(I_{L_\star})
\right]
\le C\,n^{-\beta/(2\beta+1)}.
\]
Moreover, \(I_{L_\star}\in\mathcal I_{\beta,n,0}\), the honest interval class of \cref{obj:def:honest-class}.
\end{theoremv}

\begin{remarkv}{def:frontier-handle}[Endpoint precision handle]
The endpoint precision handle comprises the following objects, for \(0<\beta\le1\), \(n\in\mathbb N\), and \(\sigma\in\mathbb R\), given the classical Legendre identities in \cref{obj:lem:classical-legendre-facts}.
\begin{itemize}
\item \textbf{(Target.)} The cutoff causal contrast
\[
\theta(P)=\mu_1(P,0)-\mu_0(P,0).
\]
\item \textbf{(Variance.)} For every \(L\in\mathbb N\), the finite derivative sum
\[
m_L=\deg K_L,\qquad
V_L(\sigma)=\frac34\sum_{j=0}^{m_L}
\frac{\sigma^{2j}}{j!}
\int_0^1\bigl(K_L^{(j)}(x)\bigr)^2\,dx,
\]
where \(K_L\) is defined in \cref{obj:def:endpoint-kernel}.
\item \textbf{(Certificate.)} The selected public radius
\[
E_\beta(n,\sigma)=\rho_{L_\star},
\]
with the radius and selected index specified in \cref{obj:def:public-selection}.
\item \textbf{(Cancellation witnesses.)} For every integer \(m\ge2\), when \(0<\sigma\le1\), and for either sign, the alternative laws
\[
b_m(\sigma)=\min\left\{1,\frac{\sigma\sqrt m}{8}\right\},
\qquad
P^\pm_{b_m(\sigma),m},
\]
defined in \cref{obj:def:alternatives,obj:def:noise-support}.
\item \textbf{(Direct witnesses.)} For every \(0<h\le1\) and either sign, the alternative laws \(P_h^{\mathrm{dir},\pm}\) defined in \cref{obj:def:direct-alternatives}.
\item \textbf{(Rate and procedures.)} The resolution rate, selected measurable estimator, and selected interval procedure
\[
r_\beta(n,\sigma),\qquad
\widetilde\theta_{\beta,n,\sigma},\qquad
\widetilde I_{\beta,n,\sigma},
\]
defined in \cref{obj:def:uniform-rate}.
\end{itemize}
\end{remarkv}

\begin{propositionv}{thm:uniform-frontier-resolution}[Uniform precision frontier]
Fix \(\beta\in(0,1]\), and write \(\nu=2\beta+1\). For \(n\ge2\) and \(\sigma\in[0,1]\), put \(h_0=n^{-1/\nu}\) and define
\[
\Psi(h,\sigma)=
\begin{cases}
0, & \sigma\le h,\\
(\sigma/h)^{2/3}-1, & \sigma^4\le h<\sigma,\\
h^{-1/2}\bigl\{1+\log(\sigma^2h^{-1/2})\bigr\}-1,
    & 0<h<\sigma^4.
\end{cases}
\]
There is a unique \(h_\star(n,\sigma)\in[h_0,1]\) satisfying
\[
\nu\log\!\left(\frac1{h_\star(n,\sigma)}\right)
+\Psi\bigl(h_\star(n,\sigma),\sigma\bigr)=\log n.
\]
Set
\[
r_\beta(n,\sigma)=h_\star(n,\sigma)^\beta,\qquad
\widetilde\theta_{\beta,n,\sigma}=\widehat\theta_{L_\star},
\qquad
\widetilde I_{\beta,n,\sigma}=I_{L_\star}.
\]

There exist constants \(c,C>0\), depending on \(\beta\) and independent of \(n\) and \(\sigma\), such that, simultaneously for every integer \(n\ge2\) and every \(\sigma\in[0,1]\),
\[
\begin{aligned}
c\,r_\beta(n,\sigma)
&\le R_\beta(n,\sigma)
\le \sup_{P\in\mathcal M_\beta(\sigma)}
   \mathbb E_{P,n}
   \left|\widetilde\theta_{\beta,n,\sigma}-\theta(P)\right|
\le E_\beta(n,\sigma)
\le C\,r_\beta(n,\sigma),\\
c\,r_\beta(n,\sigma)
&\le \mathcal L_\beta(n,\sigma)
\le \sup_{P\in\mathcal M_\beta(\sigma)}
   \mathbb E_{P,n}
   \operatorname{length}\bigl(\widetilde I_{\beta,n,\sigma}\bigr)
\le 2E_\beta(n,\sigma)
\le C\,r_\beta(n,\sigma),\\
\widetilde I_{\beta,n,\sigma}
&\in\mathcal I_{\beta,n,\sigma}.
\end{aligned}
\]
The estimation and interval comparisons range over the declared measurable estimator and honest connected-interval decision classes, including procedures with an independent uniform random seed.

The resolution has the following branches:
\begin{itemize}
\item \textbf{(Direct branch.)}
If \(\sigma\le n^{-1/\nu}\), then
\[
h_\star(n,\sigma)=n^{-1/\nu}.
\]

\item \textbf{(Intermediate branch.)}
If \(\sigma>n^{-1/\nu}\) and
\[
\log(n\sigma^{4\nu})\le \sigma^{-2}-1,
\]
then
\[
h_\star(n,\sigma)=\sigma z^{-3/2},
\qquad
z+\frac{3\nu}{2}\log z=1+\log(n\sigma^\nu),
\qquad
1<z\le\sigma^{-2},
\]
where the displayed equation uniquely determines \(z\).

\item \textbf{(Compact branch.)}
If \(\sigma>n^{-1/\nu}\) and
\[
\log(n\sigma^{4\nu})\ge \sigma^{-2}-1,
\]
then
\[
h_\star(n,\sigma)=\tau^{-2},
\qquad
2\nu\log\tau+
\tau\bigl\{1+\log(\sigma^2\tau)\bigr\}=\log n+1,
\qquad
\tau\ge\sigma^{-2},
\]
where the displayed equation uniquely determines \(\tau\).
\end{itemize}
The intermediate and compact equations agree at
\(h_\star=\sigma^4\); the direct interface is
\(h_\star=\sigma=n^{-1/\nu}\). The bounds use the same constants along every sequence \(\sigma=\sigma_n\in[0,1]\), including shrinking-noise sequences.

In particular, for \(\beta=1\) and \(\sigma=n^{-1/6}\), the intermediate branch applies for every \(n\ge2\), and
\[
r_1(n,n^{-1/6})
\asymp
\frac{n^{-1/6}}{\bigl[1+\tfrac12\log n\bigr]^{3/2}},
\]
with comparison constants independent of \(n\). Throughout, the observed side label is \(D=\mathbf1\{X\ge0\}\).
\end{propositionv}
